% Abstract (draft; numbers finalized from the data pipeline)
In sports betting, the closing odds of major bookmakers are widely
regarded as an efficient aggregation of public information, making
predictive gains over the market difficult to achieve and even harder
to verify. This paper studies forecasting under this efficiency
constraint: not whether a model can beat the market, but whether the
reliability of a forecast can be assessed before the outcome is
known, and presents a leakage-controlled evaluation of statistical,
machine-learning, and large-language-model (LLM) methods for football
match outcome forecasting, built on 11,811 matches from the five
major European leagues across seven seasons (2019/20--2025/26), with
a temporally held-out test period of 1,104 matches (the 2025/26
season to date). All features are restricted
to pre-match information, all normalization statistics are fit on the
training split only, and every reported quantity is accompanied by a
bootstrap confidence interval. \emph{The held-out test period
(2025/26 season to date) is incomplete (through 2026-02-12):
financial point estimates may shift when the season completes, so ROI
values should be read as indicative rather than definitive.} Three
findings stand out. First,
neither tree-based models nor a domain-enhanced LLM exceed the
predictive accuracy of de-vigged closing odds (54.2\% accuracy,
0.967 log loss): the market remains unbeaten, and the LLM matches
market-level probability calibration at a cost of \$0.55 for the
entire test period---an input-ablation (Table~\ref{tab:llmablation})
shows this calibration is driven by the market signal provided in
the prompt, not by independent reasoning over non-market features.
Second, permutation importance and ablations
show that market-derived features account for roughly 87\% of the
aggregate permutation importance, while referee statistics provide no measurable gain.
Third, and most importantly, forecast reliability is heterogeneous
and identifiable ex ante: an uncertainty index built from
prediction confidence and market volatility reliably stratifies
risk, with low-uncertainty matches predicted at 68--69\% accuracy
and high-uncertainty matches losing 9.6--13\% ROI when bet. In the
language of selective prediction, the resulting no-bet rule is an
abstention policy: it trades coverage against risk, and uncertainty
filtering improves risk-adjusted returns, whereas value-based
selection of high-confidence bets produces consistently negative
returns in the evaluated period, indicating that apparent model edge
is overconfidence rather than information. These results provide a
rigorous benchmark for sports forecasting research and evidence that
the decision-level objective in an efficient market is selective
prediction: abstaining when a forecast is unreliable is what connects
prediction quality to decision outcomes.
